\chapter{Graphes, chemins et matrices d’adjacence}
\section*{Exercice 34.1 -- Degrés}
La somme des degrés vaut $14$. Chaque arête contribue $2$, donc le graphe possède $7$ arêtes. Cette égalité est nécessaire ; elle ne suffit pas à elle seule pour réaliser une liste arbitraire de degrés.

\section*{Exercice 34.2 -- Adjacence}
La matrice est $M=\begin{pmatrix}0&1&0\\1&0&1\\0&1&0\end{pmatrix}$. Son carré est $\begin{pmatrix}1&0&1\\0&2&0\\1&0&1\end{pmatrix}$. Les deux nombres demandés valent $1$ : $A,B,A$ et $A,B,C$. Une marche peut répéter un sommet.

\section*{Exercice 34.3 -- Connexité}
Il n'y a aucune arête entre composantes. Chacune peut être complète : $\binom22+\binom32=1+3=4$ arêtes. La réalisation par un segment et un triangle atteint cette borne.

\section*{Exercice 34.4 -- Sommets de degré impair}
La somme de tous les degrés est paire. Les degrés pairs contribuent une somme paire ; la somme des degrés impairs doit donc être paire. Or une somme d'entiers impairs a la parité de leur nombre. Ce nombre est donc pair.

\section*{Exercice 34.5 -- Chemins et puissances}
Classer chaque marche selon son avant-dernier sommet $k$. Les $n$ premières arêtes forment une marche de $i$ à $k$, comptée par $(M^n)_{ik}$. La dernière arête de $k$ à $j$ est comptée par $M_{kj}$. Les choix se multiplient et les classes disjointes s'additionnent.

\section*{Exercice 34.6 -- Orientation}
À chaque sommet, une seule arête sortante impose la marche. Les positions se répètent tous les trois pas. Comme $2026\equiv1\pmod3$, la marche unique aboutit à $B$ ; le nombre demandé vaut $1$.
